\section*{ID7019: Full (Physics) Model — PDE System (Compact) (2)}
\paragraph{Metadata.}
PiID: 9737425626901 | RTEID: RTE:P27256.3358:T0.9278 | UUID: ebc4cd47-a71d-5cf8-93d7-cbad178bf5ae

\paragraph{Canonical Statement.}
**A. Navier–Stokes with magnetic body force (ferrohydrodynamics):** \textbackslash{}[ \textbackslash{}rho\textbackslash{}!\textbackslash{}left(\textbackslash{}frac\{\textbackslash{}partial\textbackslash{}mathbf\{u\}\}\{\textbackslash{}partial t\} + (\textbackslash{}mathbf\{u\}\textbackslash{}!\textbackslash{}cdot\textbackslash{}!\textbackslash{}nabla)\textbackslash{}mathbf\{u\}\textbackslash{}right) = -\textbackslash{}nabla p + \textbackslash{}mu \textbackslash{}nabla\textasciicircum{}2\textbackslash{}mathbf\{u\} + \textbackslash{}mathbf\{f\}\_\textbackslash{}text\{mag\} + \textbackslash{}mathbf\{f\}\_\textbackslash{}text\{ext\}, \textbackslash{}] with incompressibility \textbackslash{}(\textbackslash{}nabla\textbackslash{}cdot\textbackslash{}mathbf\{u\}=0\textbackslash{}). The magnetically induced body force may be represented (one common form) by \textbackslash{}[ \textbackslash{}mathbf\{f\}\_\textbackslash{}text\{mag\} \textbackslash{}approx \textbackslash{}mu\_0(\textbackslash{}mathbf\{M\}\textbackslash{}cdot\textbackslash{}nabla)\textbackslash{}mathbf\{H\} \textbackslash{}quad\textbackslash{}text\{(or variations used in ferrofluid literature)\}. \textbackslash{}] This term couples the applied magnetic field and local magnetization to fluid momentum and can produce normal-field (Rosensweig) instabilities on interfaces. citeturn1search0turn1search3

\paragraph{Canonical Equation.}
\[
\mathbf{f}_\text{mag} \approx \mu_0(\mathbf{M}\cdot\nabla)\mathbf{H} \quad\text{(or variations used in ferrofluid literature)}.
\]

\paragraph{Dictionary.}
A. Navier–Stokes with magnetic body force (ferrohydrodynamics): with incompressibility \textbackslash{}( =0\textbackslash{}) [ID7019]

\paragraph{Encyclopedia.}
Full (Physics) Model — PDE System (Compact) (2) is a differential equation, written f\_mag ≈ μ\_0(M·∇)H (or variations used in ferrofluid literature).. It relates the quantity to its derivatives, governing how the system evolves in space and time. It is built from ∇, f, M, H. Within the Trawin Zero Point registry it serves as one element of the framework's mathematical narrative.
